The increasing volume and complexity of graph-structured data have led to the widespread adoption of distributed graph databases in domains such as social networks, recommendation systems, knowledge graphs, and cybersecurity. By distributing graph data across multiple storage nodes, these systems can scale to billions of vertices and edges while supporting parallel query execution. However, the efficiency of a distributed graph database largely depends on the quality of graph partitioning. Poor partitioning increases the amount of inter-node communication during query execution, leading to higher latency and reduced overall system performance.

Traditional graph partitioning algorithms, including METIS \cite{MetisOG}, Kernighan--Lin, and more recent streaming approaches such as Fennel \cite{Fennel}, primarily optimize graph structure by minimizing the number of edges crossing partition boundaries while maintaining balanced partition sizes. Although minimizing the cut size is an important objective, it does not necessarily correspond to improved query performance. In many real-world applications, only a subset of graph edges is frequently traversed. Consequently, even a partitioning with a relatively small cut size may result in substantial communication overhead if frequently accessed edges are distributed across different partitions. An example of this situation is shown in Fig.~\ref{fig:intro-example}, where the colored edges represent the most frequently traversed edges. The partitioning on the left minimizes the edge cut; however, many of the cut edges are heavily traversed during query execution, resulting in a large number of Inter-Partition Traversals (IPTs). In contrast, the partitioning on the right has a bigger edge cut but significantly fewer IPTs.

\begin{figure*}[h!t]
\centering
\includegraphics[width=0.9\textwidth]{storages.png}
\caption{Example illustrating that a minimum edge cut does not necessarily minimize inter-partition traversals. Colored edges represent the most frequently traversed edges.}
\label{fig:intro-example}
\end{figure*}

To address this limitation, several workload-aware partitioning techniques have been proposed. These methods incorporate information about the actual query workload into the partitioning process, seeking to reduce the number of inter-partition traversals during query execution. Examples include Schism \cite{schism} and WAWPart \cite{wawpart}, which optimize data placement according to expected access patterns rather than relying solely on graph topology.

This paper investigates whether combining structural and workload-aware partitioning techniques can improve partition quality from the perspective of query execution. Specifically, we evaluate combinations of the Fennel streaming partitioning algorithm with the Kernighan--Lin refinement algorithm, Schism, and WAWPart. The evaluated approaches are compared using the Inter-Partition Traversal (IPT) metric, which measures the number of graph traversals crossing partition boundaries during query execution. Unlike traditional graph cut metrics, IPT directly reflects the communication overhead experienced by distributed graph databases. Furthermore, unlike execution time, IPT is independent of the implementation details and performance characteristics of a particular graph database system, making it a more objective metric for evaluating partitioning quality.

This paper presents a comparative evaluation of hybrid graph partitioning strategies that combine the Fennel streaming partitioning algorithm with refinement and workload-aware optimization techniques, including Kernighan--Lin, Schism, and WAWPart.
